\textbf{Dr. Hősi Rezső Rudolf (Eötvös Loránd University)}
\begin{itemize}
  \item \textbf{Age:} 42
  \item \textbf{Career path \& Current position:} Senior Researcher at ELTE, Department of Software Technology. PhD in Computer Vision.
  \item \textbf{Professional experience:} 18 years in high-performance computing, 3D rendering architectures, and spatial data processing.
  \item \textbf{Significant achievements:}
    \begin{enumerate}
      \item Development of the "XR-Render" framework for mobile optimization.
      \item Publication: "Efficient Point-Cloud Processing in Outdoor Environments," IEEE Transactions, 2022.
      \item Leader of the "Digital Heritage Visualization" sub-group at ELTE.
    \end{enumerate}
  \item \textbf{Distinctions:} Member of the Hungarian Academy of Sciences (MTA) Public Body.
\end{itemize}

\noindent
\textbf{Prof. Dr. Jean-Luc Soulié (Arts et Métiers)}
\begin{itemize}
  \item \textbf{Age:} 52
  \item \textbf{Career path \& Current position:} Director of the CAPLAB Platform at Arts et Métiers. Expert in 4D reality capture and XR technologies.
  \item \textbf{Professional experience:} 25+ years in VR/AR development and technology transfer from industrial sectors to cultural heritage.
  \item \textbf{Significant achievements:}
    \begin{enumerate}
      \item "Laval Virtual" Innovation Award for 4D motion capture in complex outdoor environments.
      \item Prototype: "VEP Engine" (Virtual Europe Platform) for high-fidelity 3D rendering.
      \item Standardized protocol for 4D data capture in fragile UNESCO sites.
    \end{enumerate}
  \item \textbf{Distinctions:} Member of the European XR Alliance; French National Order of Merit for Research.
\end{itemize}

\noindent
\textbf{Dr. Helen Zhang (Springer Nature)}
\begin{itemize}
  \item \textbf{Age:} 45
  \item \textbf{Career path \& Current position:} VP of Open Science Strategy at Springer Nature.
  \item \textbf{Professional experience:} Specialist in scientific publishing, Open Access revenue models, and data-sharing ecosystems for digital assets.
  \item \textbf{Significant achievements:}
    \begin{enumerate}
      \item Implementation of the "Digital Heritage + XR" Open Access publishing track.
      \item Deliverable: "Final Sustainability and Exploitation Plan" for pan-European research alliances.
      \item White Paper: "Digital Heritage Policy in the EU," presented to EU policymakers in Brussels.
    \end{enumerate}
  \item \textbf{Distinctions:} Committee Member of COPE (Committee on Publication Ethics).
\end{itemize}

\noindent
\textbf{Prof. Giulia Rossi (Sapienza Università di Roma)}
\begin{itemize}
  \item \textbf{Age:} 49
  \item \textbf{Career path \& Current position:} Full Professor of Archaeology at Sapienza. Expert in Mediterranean urban conservation.
  \item \textbf{Professional experience:} Specialist in Greco-Roman monuments and climate risk assessments for UNESCO urban sites.
  \item \textbf{Significant achievements:}
    \begin{enumerate}
      \item Development of the "Vulnerability and Tourism Impact Index" for historical centers.
      \item Co-author: "Climate Risks in Mediterranean Heritage," Nature Heritage, 2022.
      \item Lead archaeologist for the digital documentation of the Roman Forum ruins.
    \end{enumerate}
  \item \textbf{Distinctions:} UNESCO Chair on Urban Heritage Preservation.
\end{itemize}

\noindent
\textbf{Prof. Dr. Dimitri Papadopoulos (National and Kapodistrian University of Athens)}
\begin{itemize}
  \item \textbf{Age:} 55
  \item \textbf{Career path \& Current position:} Full Professor of Byzantine Archaeology at NKUA.
  \item \textbf{Professional experience:} Specialist in long-term site monitoring and climatic risk analysis for open-air Mediterranean ruins.
  \item \textbf{Significant achievements:}
    \begin{enumerate}
      \item Lead investigator for the "Climate-Heritage" EU-sensor network for ancient sites.
      \item Publication: "Erosion Patterns in Eastern Mediterranean Marble Structures," Archaeometry Today, 2021.
      \item Developer of the "Aegis-Risk" database for Byzantine fortification monitoring.
    \end{enumerate}
  \item \textbf{Distinctions:} Member of the International Council on Monuments and Sites (ICOMOS).
\end{itemize}

\noindent
\textbf{Dr. Claire Lefebvre (Sorbonne Université)}
\begin{itemize}
  \item \textbf{Age:} 41
  \item \textbf{Career path \& Current position:} Senior Researcher at the Institut d'art et d'archéologie, Sorbonne Université.
  \item \textbf{Professional experience:} Expert in medieval and early modern architectural history, specifically focusing on Atlantic coastal fortifications.
  \item \textbf{Significant achievements:}
    \begin{enumerate}
      \item Coordination of the "Atlantic Wall" digital documentation project (2018–2022).
      \item Publication: "Coastal Erosion and the Loss of Medieval Heritage," European Journal of Architecture, 2024.
      \item Development of the "Fortress-Data" documentation protocol for military heritage.
    \end{enumerate}
  \item \textbf{Distinctions:} Recipient of the Ordre des Arts et des Lettres for cultural heritage research.
\end{itemize}
  
  